\chapter{Introduzione}

In tempi recenti, il Digital Twin si è affermato come una delle tecnologie più trasformative dell'Industria 4.0.
\noindent Inizialmente concepito in ambito aerospaziale allo scopo di monitorare lo stato dei veicoli spaziali, il gemello digitale è diventato rapidamente una tecnologia trasversale in tutti i campi aperti al progresso tecnologico. Chiave del suo successo è la fusione di modelli matematici e dati reali, il gemello digitale permette la realizzazione di strumenti che spaziano dalla predizione di guasti alla coordinazione di complessi sistemi e processi.

In contemporanea alla sua diffusione in ambito industriale, al fine di monitorare lo stato degli impianti e macchine di produzione, si è assistita a una rapida crescita dell'interesse accademico alla tecnologia, tale da scaturire una rapida proliferazione di interessanti articoli scientifici estesi non più solo in campi industriali, ma anche in settore primario e servizi.

Diverse ricerche sui digital twin hanno esplorato soluzioni nell'ambito delle \textit{smart cities} \cite{dt_smart_cities} e l'integrazione di gemelli digitali con dispositivi IoT nella pianificazione urbana; in agricoltura sono state studiate applicazioni per il monitoraggio delle colture \cite{dt_agricolture}; nelle \textit{smart grids} \cite{smartgerudo} nel contesto della distribuzione energetica in combinazione all'uso di intelligenza artificiale. Come è facile evincere, la ricerca si propaga nella direzione di sistemi sempre complessi di gemelli digitali.

L'integrazione di gemelli digitali in sistemi multi-agente è di chiave importanza in molti dei contesti citati: nelle smart cities, veicoli a guida autonoma possono coordinare percorsi al fine di ridurre congestioni; in agricoltura gemelli digitali di robot possono collaborare nel campionamento delle colture al fine di individuare eventuali fitopatie nelle piante.

Tali sistemi multi-agenti richiedono opportune infrastrutture software, scalabili e robuste, al fine di supportare cooperazioni su larga scala.


Il presente lavoro di tesi si occuperà di sviluppare una implementazione di un framework per la realizzazione di sistemi di gemelli digitali multi-agente distribuito, fondando la sua base teorica sull'architettura presentata dall'articolo scientifico \textit{Advancing Robotic Systems with Distributed Multi-Agent Digital Twins: A Scalable and Adaptive Framework} \cite{Russo2025}.

Dai principi e obiettivi fissati nella pubblicazione di riferimento, è stata realizata una suite di software in grado di collegare dispositivi embedded e simulazione nel motore \textbf{Godot Engine}, implementare gemelli digitali con questo collegamento, e orchestrare un ambiente virtuale condiviso tramite le API per il networking di Godot stesso.

\bigbreak 

%I Digital Twin (repliche virtuali di sistemi fisici) sono diventati strumenti fondamentali nella robotica, nell'industria 4.0, nella sanità e nelle smart city. Permettono di simulare scenari, prevedere guasti e ottimizzare le prestazioni in tempo reale senza rischi, riducendo costi e tempi di fermo. L'integrazione con l'Intelligenza Artificiale (IA) sta portando a sistemi sempre più autonomi e intelligent\section{Caso di studio}
 

\subsubsection*{Struttura della tesi}
Nei successivi capitoli verranno trattati i seguenti argomenti:
\begin{itemize}
  \item Il \textbf{Capitolo 2} descrive il modello del DMDTE.

  \item Il \textbf{Capitolo 3} introduce le tecnologie software e hardware utilizzate.
  \item Il \textbf{Capitolo 4} descrive nel dettaglio l'implementazione di DMDTE, descrivendone architettura, implementazione e uso.
  \item Il \textbf{Capitolo 5} discute delle varie sperimentazioni eseguite con il framework. 
\end{itemize}